\documentclass[10pt,a4paper]{amsart}
\usepackage[margin=19mm]{geometry}
\usepackage{amsmath,amssymb,mathtools}
\usepackage[scr=rsfs,cal=euler]{mathalfa}
\usepackage{xfrac}
\usepackage[unicode,hidelinks,bookmarks=false]{hyperref}
\newcommand{\ie}{i.e.\ }
\newcommand{\cf}{cf.\ }
\newcommand{\resp}{respectively\ }
\newcommand{\Subsection}{\subsection}
\providecommand{\nobreakdash}{\nobreak\hbox{-}}
\setlength{\emergencystretch}{2em}
\multlinegap0pt
\usepackage[shortlabels]{enumitem}
\usepackage{amssymb}
\newtheorem{lem}{Lemma}
\newtheorem{prop}{Proposition}
\usepackage{cleveref}
\crefname{lem}{Lemma}{Lemmas}
\crefname{prop}{Proposition}{Propositions}
\title{Conditional means, vector pricings,\\ amenability and fixed points in cones}
\author{Nicolas Monod}
\date{Source: arXiv:2512.13829v4 (CC BY 4.0)}
\begin{document}
\maketitle

\noindent\emph{Source and licence.} Conditional means, vector pricings,\\ amenability and fixed points in cones. Version arXiv:2512.13829v4 (CC BY 4.0), distributed by arXiv under the Creative Commons Attribution 4.0 International licence, http://creativecommons.org/licenses/by/4.0/, as recorded in the arXiv licence metadata for this exact version and shown on the abstract page https://arxiv.org/abs/2512.13829v4. Full licence notice: \texttt{licenses/arxiv-2512.13829-CC-BY-4.0.txt}.\par\smallskip

\noindent\textit{Editorial scope.} Counted unit: the article's prop prop:VP:CM:equiv (source0.tex lines 596--604). The article's complete original proof of it is delivered with it (source0.tex lines 606--641, one \texttt{proof} environment of 36 lines). No mathematics is written, repaired or re-derived: the statement and the proof are the article's own lines, in the article's own notation, and no retained line is substituted or re-flowed.  Retained uncounted as the source-local context the proof needs: lines 219--224; lines 540--542; lines 567--588.  Nothing was omitted from any retained span, so the package carries no omission record.  No retained line contains a citation, so the wrapper carries no bibliography and no bibliography entry.  No retained line contains a figure environment, an image-inclusion command or drawing code, so no image is delivered and the package declares no asset dependency.  Editorial formatting only, in addition: an AMS \texttt{amsart} preamble that restates the article's own declarations and macros used by the retained lines, namely the environments \texttt{lem}, \texttt{prop} on separate counters, together with the standard packages those lines need.  The retained environments print in this document as Lemma 1, Proposition 1 in order of appearance, whereas the article prints the same items with the numbering of its own full document.  The complete native source of the article as distributed in the arXiv e-print for this exact version is bundled unchanged as \texttt{sources/191058/source0.tex}.
\smallskip
\begin{enumerate}[left=2\parindent, labelsep=\parindent, topsep=3pt, itemsep=3pt, label=\normalfont(VP\arabic*)]
\item $r(u+v, w) = r(u, w) + r(v, w)$ when $w\neq 0$. \label{pt:VP:add}
\item $ r(u,w) = r(u, v) \, r(v, w)$ when defined. \label{pt:VP:cocycle}
\item $r(u,u)=1$.\label{pt:VP:one}
\end{enumerate}

\begin{lem}\label{lem:no-0}
If a map $r$ defined on $V^+ \times ( V^+ \smallsetminus \{0\})$ satisfies the conditions for vector pricings on its domain of definition, then it extends (uniquely) to a vector pricing by setting $r(0,0)=1$ and $r(u, 0)=+\infty$ when $u\neq 0$.\qed
\end{lem}

\section{Pricings vs means vs partial functionals}\label{sec:PPP}%%%%%%%%%%%%%%%%%%%%%%%%%%%%%%%%%%%%%%%%%%%%%%%%%%%%%%%%%%%%%%%%%%
%%%%%%%%%%%%%%%%%%%%%%%%%%%%%%%%%%%%%%%%%%%%%%%%%%%%%%%%%%%%%%%%%%%%%%%%%%%%%%%%%%%%
In this section, we establish dictionaries between three languages for the results of this article.
\subsection{Conditional means}\label{sec:CM}%%%%%%%%%%%%%%%%%%%%%%%%%%%%%%%%%%%%%%%%%%%%%%%%%%%%%%%%%%%%%%%%%%
\begin{flushright}
\begin{minipage}[t]{0.85\linewidth}\itshape
One of the common denominators I have found is that expectations rise above that which is expected.
\end{minipage}
\end{flushright}
\begin{flushright}
\begin{minipage}[t]{0.75\linewidth}\small
G.W. Bush, Los Angeles, 27 September 2000\\
(c.f. Congressional Record, Vol.~146 (2000), Part 17)
\end{minipage}
\end{flushright}
% 
We extend the definition of conditional means from $\ell^\infty$ to the general case. Let $V$ be an ordered vector space. A \textbf{conditional mean} on $V$ is a function $P$ with values in $[0,1]$ defined on all pairs of vectors $0\leq u\leq v \neq 0$ such that the following hold.% (whenever the variables are in the domain of definition).
\begin{enumerate}[left=2\parindent, labelsep=\parindent, topsep=3pt, itemsep=3pt, label=\normalfont(CM\arabic*)]
\item $P(u+v | w) = P(u| w)+P(v | w)$ when $u+v\leq w\neq 0$ with $u,v\geq 0$. \label{pt:CM:add}
\item $P(u|w) = P(u|v)\, P(v|w)$ when $0\leq u \leq v \leq w$ with $v\neq 0$.\label{pt:CM:trans}
\item $P(u|u)=1$ when $u\gneqq 0$.\label{pt:CM:one}
\end{enumerate}

\begin{prop}\label{prop:VP:CM:equiv}
Let $P$ be a conditional mean on the ordered vector space $V$. Then the expression
\begin{equation*}
r_P(u,v) = \frac{P(u|u+v)}{P(v|u+v)}\kern3mm (u,v\in V^+, v\neq 0)
\end{equation*}
defines a vector pricing on $V$.

Moreover, the assignments $r\mapsto P_r$ and $P\mapsto r_P$ are mutually inverse.
\end{prop}

\begin{proof}
We only consider $r_P(u,v)$ when $v\neq 0$, see \Cref{lem:no-0}. Since $P(\cdot| u+v)$ is additive with $P(u+v| u+v)=1$, it cannot happen that both numerator and denominator vanish. Thus, $r_P(u,v)$ is well-defined in $[0, +\infty]$.

Observe that whenever $x\geq u+v$ is such that $P(u+v|x) \neq 0$, \ref{pt:CM:trans} implies
\begin{equation*}
\frac{P(u|x)}{P(v|x)} = \frac{P(u|u+v) P(u+v|x)}{P(v|u+v) P(u+v|x)} = \frac{P(u|u+v)}{P(v|u+v)} =r_P(u,v).
\end{equation*}
In order to check \ref{pt:VP:add}, consider $u,v,w\in V^+$ with $w\neq 0$. If both
\begin{equation*}
P(u+w|u+v+w) \neq 0 \kern3mm \text{and} \kern3mm  P(v+w|u+v+w)\neq 0,
\end{equation*}
then the above observation with $x=u+v+w$ allows us to verify
\begin{multline*}
r_P(u, w) + r_P(v, w) = \frac{P(u|u+v+w)}{P(w|u+v+w)} +  \frac{P(v|u+v+w)}{P(w|u+v+w)} =  \\
= \frac{P(u+v|u+v+w)}{P(w|u+v+w)} = r_P(u+v, w).
\end{multline*}
Otherwise, without loss of generality $P(u+w|u+v+w)=0$ which implies $P(w|u+v+w)=0$ and thus $r_P(u+v, w)=+\infty$. On the other hand, $P(u+w|u+v+w)=0$ also implies $P(v|u+v+w)=1$ and hence also $P(v+w|u+v+w)=1$. Combining this with \ref{pt:CM:trans} applied as
\begin{equation*}
P(w|v+w) \, P(v+w|u+v+w) = P(w|u+v+w)=0
\end{equation*}
implies $P(w|v+w)=0$ and hence $r_P(v, w)=+\infty$, so that \ref{pt:VP:add} holds.

The verification of \ref{pt:VP:cocycle}, namely
\begin{equation*}
\frac{P(u|u+v)}{P(v|u+v)} \, \frac{P(v|v+w)}{P(w|v+w)} =\frac{P(u|u+w)}{P(w|u+w)},
\end{equation*}
is entirely similar, considering again $x=u+v+w$. Namely, if none of the pairwise sums of $u$, $v$ and $w$ are zeros of $P(\cdot|u+v+w)$, then the identity follows from the above observation. Otherwise, two among $u$, $v$ and $w$ are zeros of $P(\cdot|u+v+w)$. If this is the case of $u$  and $v$, then $P(w|u+v+w)=1$, which implies $P(u+w|u+v+w)=1$ and hence using \ref{pt:CM:trans} for
\begin{equation*}
P(u|u+w) \, P(u+w|u+v+w) = P(u|u+v+w)=0
\end{equation*}
we deduce $P(u|u+w)=0$. Likewise, $P(v|v+w)=0$; now the desired relation reads $0=0$. 

In case $v$ and $w$ are zeros of $P(\cdot|u+v+w)$, we deduce in the same way that the relation becomes $+\infty=+\infty$ and finally in case $u$ and $w$ are zeros of $P(\cdot|u+v+w)$, the left hand side is the undefined form $0\cdot+\infty$.

\ref{pt:VP:one} holds trivially and the remaining statement about $r\mapsto P_r$ and $P\mapsto r_P$ is a direct verification.
\end{proof}

\end{document}
